\section{The problem and the experimental object}\label{sec:experiments}
The geometry relevant to prediction is the geometry of information that an experiment has actually acquired. It is not, in general, the geometry of the parameter space, of all formally reachable posteriors, or of an arbitrarily selected reference measure. A second distinction is equally important. A finite list of accurate terminal predictions does not itself give a causal implementation of that list: the reports needed to select a representative may already have been discarded.

Our main result, Theorem~\ref{thm:main}, connects these two issues. It is a sufficient raw-kernel construction from a static global cover, block two-point stability, and a cover-preserved relation which admits a robust envelope bound. Unlike a one-step contraction theorem, it allows every individual transition in a block to have contraction coefficient one. Unlike a blockwise coding prescription, its encoder retains only one finite label after \emph{each} physical report. The proof compares its perturbed trajectory with an exact shadow started at the beginning of a block; a relation on candidate states, independent of the label budget, controls the perturbations inside the block. Actual acquired small-ball mass then gives the converse in the same squared task.

The distinction has a concrete consequence. Sparse transition matrices need not mix in one step, and continuous Gaussian observations have unbounded log-likelihood ratios. Nevertheless, a primitive support pattern and a positive-probability bounded observation block suffice. A mass-preserving inward compander prevents intermediate rounding from destroying the positivity accumulated over that block. The resulting hard-label law has exponent determined by a positive submeasure of the actual posterior distribution. Theorem~\ref{thm:hmm} proves this for an arbitrary fixed finite dimension, with time-varying sparse matrices and state-dependent missing observations. Theorem~\ref{thm:singular} verifies the same general mechanism on an entirely different, singular predictive domain. There the actual acquired law is finite atomic; a physical-oracle comparison, not a substitution of a Cantor law for that atomic law, supplies the geometric converse.

The general theorem is followed by exact acquisition identities, a stopped drift estimate, and a composition theorem for fully specified causal transducers. Uniform deterministic-time risk is distinguished from stopped risk, and a complete report-law comparison is distinguished from a prediction bound. The structural hypotheses are not asserted to be necessary for every finite-memory predictable experiment. In particular, exact copying can preserve a code without any strict contraction. This boundary does not reduce the experimental problem; it identifies which mechanism the present theorem proves.

\subsection{Prepared kernels and legal strategies}
All measurable spaces below are standard Borel. A parameter-indexed experiment consists of hidden spaces $X_t$, action spaces $A_t$, report spaces $Y_{t+1}$, a preparation $\mu_\lambda$, and measurable probability kernels
\begin{equation}\label{eq:raw}
 K_t^\lambda(dx',dy,dc\mid x,a).
\end{equation}
The cost $c$ takes values in a declared finite product of nonnegative extended real spaces and includes elapsed physical time. Positivity means nonnegativity and normalization; no strict lower bound is imposed at this level. Failures, missing observations and Dirac transitions are legitimate reports. A visible history records actions, reports and public costs. A continuation interface records available actions, remaining resources and any clock which changes the legal future.

A strategy is a measurable kernel $\beta_t(da\mid h_t)$ supported on the available actions. It may use independent randomization but cannot read $\lambda$, an unobserved hidden state, or a future report. A stopping rule is a stopping time for this visible filtration, with the declared stop action and terminal test. Kernel integration defines the actual path law. For the Bayesian prediction statements, fix a preparation (or a prior on $(\lambda,X_0)$) and a legal exploration rule \emph{independently of the encoder}. A posterior for this preparation is not automatically a parameter-uniform sufficient statistic. The latter quantifiers occur explicitly in Section~\ref{sec:morphisms}.

A persistent $M$-label encoder has measurable updates $I_t=q_t(I_{t-1},a_t,y_t)$ with at most $M$ possible labels at each cut. Time-dependent read-only programs are allowed only with their description and clock accounted for. Randomized updates are allowed unless stated otherwise. No earlier report can be reread. A checkpoint encoder inspects the actual terminal history once and retains at most $M$ labels. Neither model includes a free real posterior, unbounded random seed containing data, or hidden observation tape. A supplied controller's state is charged separately from the prediction label.

\subsection{Executable future tests and a measurable quotient}
An executable test is a legal finite continuation followed by a bounded function of its actual reports. Histories are equivalent when they have the same continuation interface and all conditional test expectations agree. This definition concerns an operational class of tests, not arbitrary functions of a hidden state. The next certificate supplies a measurable realization; applications may instead verify their quotient directly.

\begin{proposition}[Continuous-instrument quotient]\label{prop:quotient}
Let $\mathcal B$ be compact metrizable. Let $\mathcal V\subset C(\mathcal B)$ be a separable closed linear space generated by continuous executable-test evaluations and interface coordinates, containing constants. A countable dense family determines all declared tests. For each action $a$, let the report domain be Polish with reference measure $m_a$, and let $g_a(\pi,y)$ be a jointly continuous nonnegative density. On $D_a=\{g_a>0\}$ suppose the normalized belief update $B_a$ is continuous. For every $f\in\mathcal V$ and report $y$, require the instrument evaluation $g_a(\pi,y)f(B_a(\pi,y))$, extended by its instrument value at density zero, to be continuous in $\pi$ and in $\mathcal V$. Require joint continuity of these extended evaluations where continuity in the report is used. Then the evaluation image $S$ is compact metrizable, its fibers are predictive equivalence, and the report laws and normalized updates descend to $S$. The descended update is continuous on the open positive-density domain in $S\times Y$.

If a standard-Borel statistic measurably determines every coordinate evaluation, it measurably determines $S$. The assertion may be imposed on one common full-measure history set for an almost-sure version.
\end{proposition}
\begin{proof}
Normalize the determining evaluations to lie in $[-1,1]$ and let $E:\mathcal B\to[-1,1]^{\N}$ be their product. Its image is compact and hence Borel. Agreement on its coordinates implies agreement on $\mathcal V$ by uniform approximation and on the declared tests by the determining hypothesis. Taking $f=1$ shows that equal fibers have the same density at every report. Dividing the equal instrument evaluations by that positive common density shows that $E(B_a(\pi,y))$ is constant on the corresponding fibers.

The continuous surjection $E$ is closed, with compact fibers. The product $E\times\mathrm{id}_Y$ is a closed quotient map: given a closed set and a convergent image point, compactness of the belief fiber supplies a convergent subnet witnessing closedness locally. Equivalently one can use the finite subcover characterization of a closed map with compact fibers. The joint instrument coordinates therefore descend continuously. Their ratios by the positive descended density are continuous on the open set where it is positive; these ratios are precisely the coordinates of the descended update. No claim is made about normalization at a null report. A jointly Borel null-report extension, when needed by a numerical representative, is additional data.

For minimality let $u_n$ be measurable coordinate factors through the given statistic. The product $u=(u_n)$ is measurable. Replace $u$ by a fixed state outside the Borel set $S$ to obtain an $S$-valued factor. Countably many almost-sure equalities can be imposed on one common full-measure set. Thus the conclusion does not rely on measurability of an unspecified set-theoretic quotient.
\end{proof}

The metric used for stability need only generate the relevant Borel structure, not the compact quotient topology. Log-ratio distance on the open probability simplex is an important example. Report versions used to update a wrong representative must be defined jointly, rather than selected separately almost surely for each state.

\subsection{One scored task}
Let $Y$ be a bounded scalar terminal report, or let an independently selected finite probe index specify one of finitely many bounded scalar reports. In the latter case use the weighted Euclidean norm of their separately executable squared losses. No joint counterfactual measurement is assumed. Put
\begin{equation}\label{eq:task}
 \begin{gathered}
 P_T=\E[Y\mid\cH_T],\quad \nu_T=\law(P_T),\quad B_T=\E\norm{Y-P_T}^2,\\
 q_M(\nu)=\inf_{\#C\le M}\int\dist(p,C)^2\nu(dp).
 \end{gathered}
\end{equation}
The infimal checkpoint and online risks are $R_\cp(T,M)$ and $R_\on(T,M)$. The law $\nu_T$ is always the pushforward of the actual preparation and exploration. It may be atomic, mixed, singular, transient, or supported on several strata. Its role in the theorem is probabilistic rather than formal-dimensional.
